% ============================================================
%  SASTO Project Notebook / Lab Logbook
%  Eric Hou — SYNOPSYS 2026
%  Engineering Category
%
%  Compile:  pdflatex project_notebook.tex  (run twice)
%
%  All images referenced from notebook_images/ directory.
%  Clean figures generated by generate_notebook_figures.py
% ============================================================

\documentclass[11pt, letterpaper]{article}

\usepackage[
  margin=1in,
  top=0.9in,
  bottom=0.9in
]{geometry}

\usepackage[T1]{fontenc}
\usepackage[utf8]{inputenc}
\usepackage{lmodern}
\usepackage{microtype}
\usepackage{graphicx}
\usepackage{xcolor}
\usepackage{enumitem}
\usepackage{booktabs}
\usepackage{tabularx}
\usepackage{fancyhdr}
\usepackage{titlesec}
\usepackage{hyperref}
\usepackage{float}
\usepackage{parskip}
\usepackage{tcolorbox}
\usepackage{amsmath}
\usepackage{longtable}

\hypersetup{
  colorlinks=true,
  linkcolor=blue!60!black,
  urlcolor=blue!60!black
}

% ── Colours ──────────────────────────────────────────────────
\definecolor{entryblue}{HTML}{1A5276}
\definecolor{failred}{HTML}{C0392B}
\definecolor{fixgreen}{HTML}{1E8449}
\definecolor{notegray}{HTML}{5D6D7E}
\definecolor{bglight}{HTML}{F2F4F4}

% ── Entry environment ────────────────────────────────────────
\newcounter{entrynum}
\newcommand{\entry}[2]{%
  \refstepcounter{entrynum}%
  \vspace{6mm}%
  \noindent\begin{tcolorbox}[
    colback=bglight, colframe=entryblue,
    boxrule=1pt, arc=3pt,
    title={\textbf{Entry \theentrynum}\hfill\textbf{#1}},
    fonttitle=\large\color{white},
    colbacktitle=entryblue,
    left=3mm, right=3mm, top=2mm, bottom=2mm
  ]
  {\large\bfseries\color{entryblue} #2}
  \end{tcolorbox}
  \vspace{-1mm}%
}

% Field labels
\newcommand{\objective}{\par\textbf{\color{entryblue}Objective:} }
\newcommand{\experiment}{\par\textbf{\color{entryblue}Experiment / Method:} }
\newcommand{\result}{\par\textbf{\color{entryblue}Result:} }
\newcommand{\failure}{\par\textbf{\color{failred}Failure / Issue:} }
\newcommand{\fix}{\par\textbf{\color{fixgreen}Design Decision / Fix:} }
\newcommand{\nextstep}{\par\textbf{\color{notegray}Next Step:} }
\newcommand{\obs}{\par\textbf{\color{entryblue}Observations:} }

% Crossed-out text for "erased" notes
\newcommand{\crossed}[1]{\textcolor{red!50!black}{\sout{#1}}}
\usepackage[normalem]{ulem}

% ── Image command ────────────────────────────────────────────
\graphicspath{{notebook_images/}}
\newcommand{\nbimg}[2][0.85\linewidth]{%
  \begin{center}
    \includegraphics[width=#1,keepaspectratio]{#2}
  \end{center}
}
\newcommand{\nbimgcap}[3][0.85\linewidth]{%
  \begin{figure}[H]
    \centering
    \includegraphics[width=#1,keepaspectratio]{#2}
    \caption{\small\itshape #3}
  \end{figure}
}

% ── Page style ───────────────────────────────────────────────
\pagestyle{fancy}
\fancyhf{}
\fancyhead[L]{\small\color{notegray} SASTO Project Notebook — Eric Hou}
\fancyhead[R]{\small\color{notegray} SYNOPSYS 2026}
\fancyfoot[C]{\small\thepage}
\renewcommand{\headrulewidth}{0.4pt}

\setlength{\parindent}{0pt}
\setlength{\parskip}{4pt}

% ============================================================
\begin{document}

% ── TITLE PAGE ───────────────────────────────────────────────
\begin{titlepage}
  \centering
  \vspace*{3cm}
  {\Huge\bfseries\color{entryblue}
    Project Notebook}\\[6pt]
  {\LARGE\bfseries
    Additive Manufacturing: Harnessing FEA to Optimize Material Efficiency}\\[1.5cm]
  {\Large Eric Hou}\\[0.5cm]
  {\large Los Gatos, California}\\[2.5cm]
  {\large\color{notegray}
    Dates covered: November 2025 -- March 2026\\[4pt]
    Total entries: 39\\[4pt]
    This notebook documents the research process,\\
    including failures, abandoned approaches, debugging,\\
    and iterative improvements.}
  \vfill
  {\small\itshape
    \textbf{Note:} This is a reconstructed digital notebook, assembled
    from dated source files, commit logs, error logs, and notes
    written during the project. It reflects the actual sequence of
    work as accurately as possible, but was typeset after the fact.
    Mistakes are crossed out rather than deleted.}
\end{titlepage}

% ── TABLE OF CONTENTS ────────────────────────────────────────
\tableofcontents
\newpage

% ════════════════════════════════════════════════════════════
\section{Phase 1: Idea Formation \& Early Exploration}
% ════════════════════════════════════════════════════════════

% ──────────────────────────────────────────────────────────────
\entry{November 14, 2025}{Initial Project Concept}

\objective
Define the core research question. I want to explore whether AI can reduce
material waste in construction --- specifically concrete, which accounts for
$\sim$8\% of global CO$_2$ emissions.

\experiment
Literature review of topology optimization (Bendsoe \& Sigmund 2003),
3D-printed concrete buildings (Buswell et al.\ 2018, ICON, COBOD), and
neural network surrogates for FEA (Nie 2021, Banga 2018). Read ASCE 7-22
load specification.

\obs
\begin{itemize}[nosep]
  \item Classical SIMP requires 100--1000 FEA solves per design --- too slow
    for building-scale meshes.
  \item 3D concrete printers can print arbitrary wall profiles at no extra cost.
  \item No existing work combines surrogate FEA with topology optimization
    \emph{and} printability enforcement for full buildings.
\end{itemize}

\nextstep
Find a dataset of building geometries to train a surrogate model on.

% ──────────────────────────────────────────────────────────────
\entry{December 2, 2025}{Web Crawler for Floor Plans}

\objective
Collect floor plan images from architectural websites to generate 3D building
models.

\experiment
Built a web crawler (\texttt{crawler/crawler.py}) targeting
architecturaldesigns.com. Used OpenAI CLIP (ViT-B/32) to classify scraped
images as floor plans vs.\ exterior renders. Set confidence threshold at 0.85.
Added perceptual hashing for deduplication.

\result
Collected $\sim$1,200 floor plan images after running the crawler overnight.
Classification threshold 0.85 seemed reasonable but I did not do a full
manual audit --- spot-checking 50 images showed mostly real floor plans.

\failure
Many images were low-resolution marketing renders, not dimensioned
architectural drawings. The images had text overlays, watermarks, and
inconsistent scales. I had no way to convert these into accurate 3D geometry
with known dimensions.

\nextstep
Investigate Plan2Scene pipeline for floor-plan-to-3D conversion.

% ──────────────────────────────────────────────────────────────
\entry{December 10, 2025}{Plan2Scene Pipeline --- First Attempt}

\objective
Use the Plan2Scene approach to convert crawled floor plan images into 3D
volumetric building models.

\experiment
Set up the Plan2Scene codebase. Built a pipeline
(\texttt{data\_model/complete\_pipeline.py}) that:
\begin{enumerate}[nosep]
  \item Detects rooms and walls from the floor plan image
  \item Extrudes walls to a fixed height
  \item Adds roof surface
\end{enumerate}

\result
Got basic 3D models from a few clean floor plans.

\failure
\begin{itemize}[nosep]
  \item Image-based detection was unreliable --- walls were frequently missed
    or misaligned.
  \item No dimensional accuracy. Wall thicknesses were arbitrary since the
    source images had no scale bar.
  \item Room detection failed on complex multi-story plans.
  \item Generated meshes were \textbf{not watertight} --- full of holes and
    self-intersections.
\end{itemize}

\fix
\sout{Try improving the image segmentation model.} After two weeks of
debugging, I concluded this approach was fundamentally limited. Image-based
geometry recovery cannot produce the dimensional accuracy needed for structural
analysis. I need a dataset with actual 3D coordinates.

\nextstep
Search for wireframe or CAD-based building datasets.

% ──────────────────────────────────────────────────────────────
\entry{January 4, 2026}{Pivot to 3DWire Dataset}

\objective
Find a proper 3D building dataset with topological connectivity information.

\experiment
Literature search found the \textbf{3DWire dataset} (Lin et al.\ 2024, KAUST) ---
14,293 building wireframe skeletons with vertex coordinates and edge connectivity.
Downloaded and explored the data format.

\result
Each wireframe is a JSON file with 3D vertex positions and edge lists encoding
walls, floors, and roof ridges. This looks like exactly what I need --- but I
still have to figure out how to ``inflate'' wireframe edges into solid volumes.

\obs
\begin{itemize}[nosep]
  \item 14,293 wireframes --- much larger than I expected.
  \item Diverse building types: rectangular, L-shaped, multi-gabled, etc.
  \item No volumetric geometry yet --- just wireframe edges. I need to
    ``inflate'' these into solid walls, floors, and roofs.
\end{itemize}

\fix
Abandoned the web crawler + Plan2Scene approach entirely. The 3DWire dataset is
a much better foundation. This was a \textbf{major pivot} --- about 5 weeks of
work on the crawler was effectively wasted, but it taught me what I actually
needed from the data.

\nextstep
Build a pipeline to convert wireframe edges into volumetric solid geometry.


% ════════════════════════════════════════════════════════════
\section{Phase 2: Dataset Generation \& Geometry Pipeline}
% ════════════════════════════════════════════════════════════

% ──────────────────────────────────────────────────────────────
\entry{January 10, 2026}{Wireframe-to-Volume Extrusion --- v1}

\objective
Write code to convert 3DWire wireframe edges into solid volumetric geometry
(walls, floors, roofs) suitable for FEA meshing.

\experiment
Built \texttt{wireframe\_to\_volume.py} (eventually grew to 1,684 lines).
First version: extrude each wall edge into a rectangular prism with fixed
thickness. Floor = horizontal slab at $z=0$. Roof = flat slab at max $z$.

\result
I successfully generated basic box-like houses for simple rectangular wireframes. The pipeline ran end-to-end for the first time.

\failure
\begin{itemize}[nosep]
  \item Walls overlapped at corners --- the Boolean union to merge them
    produced \textbf{self-intersecting geometry}.
  \item Roof was just a flat slab, not pitched. Looks unrealistic and
    structurally wrong.
  \item Wall thickness was the same everywhere --- exterior thick, interior
    thin would be more realistic.
\end{itemize}

\nextstep
Implement proper Boolean fusion and pitched roof generation.

% ──────────────────────────────────────────────────────────────
\entry{January 18, 2026}{Boolean Fusion Failures}

\objective
Merge overlapping wall/floor/roof solids into a single watertight body using
Boolean union.

\experiment
Tried using \texttt{trimesh.boolean} with the \texttt{blender} engine for
Boolean operations. Also tried \texttt{pymesh} and direct CSG.

\result
I had partial success --- simple houses merged correctly, but more complex geometries consistently broke.

\failure
\textbf{This was one of the most frustrating debugging episodes of the
entire project.}
\begin{itemize}[nosep]
  \item Boolean operations failed on $\sim$40\% of wireframes.
  \item Failure modes: self-intersecting faces, non-manifold edges,
    degenerate triangles, missing faces.
  \item blender engine was slow (8--15s per house) and crashed on complex
    geometry.
  \item trimesh kept returning empty meshes for some inputs.
\end{itemize}

\fix
Switched to \texttt{manifold3d} library which handles Boolean operations
much more robustly. Added a pre-fusion mesh cleanup step: remove degenerate
faces, snap vertices to grid, repair non-manifold edges. This reduced failure
rate from $\sim$40\% to $\sim$5\%.

\nextstep
Implement pitched roof generation using ConvexHull.

% ──────────────────────────────────────────────────────────────
\entry{January 22, 2026}{Pitched Roof and Part Labels}

\objective
Generate proper pitched roofs and assign structural part labels to each region
of the building (exterior wall, interior wall, roof, floor).

\experiment
\begin{itemize}[nosep]
  \item Roof: compute ConvexHull of upper wireframe vertices with apex raised
    by a pitch factor. Extrude hull inward to create solid roof geometry.
  \item Part labels: assign integer codes (1=exterior wall, 2=interior wall,
    3=roof, 4=floor) based on geometric position relative to building envelope.
  \item Scale-safe thickness fractions: exterior $\leq$3.0\% of building width,
    interior $\leq$2.0\%, roof $\leq$1.5\%.
\end{itemize}

\result
I now had much more realistic building geometry with proper pitched roofs. Part labels were correctly assigned. I generated a test set of 50 houses for visual inspection.

\nbimgcap[0.75\linewidth]{fig_voxel_house.png}{Voxel representation of a
generated house showing part-label coloring: blue = exterior walls, orange-red =
interior walls, green = roof, tan = floor slab.}

\obs
Part labels will be important later --- I can use them to allow different
optimization aggressiveness for different structural elements (e.g., thin
interior walls more than exterior).

\nextstep
Set up the FEA pipeline to simulate these structures.

% ──────────────────────────────────────────────────────────────
\entry{January 28, 2026}{FEA Pipeline --- First Attempts}

\objective
Run finite element analysis on the generated volumetric buildings to compute
stress, displacement, and compliance.

\experiment
Set up SfePy (Simple Finite Elements in Python) for static linear elastic
analysis. Material: concrete ($E = 25$\,GPa, $\nu = 0.20$,
$\rho = 2400$\,kg/m$^3$). Load cases per ASCE 7-22: dead load, live load,
wind combinations.

\result
I got the first FEA solve to complete without crashing. The stress contours
\emph{looked} physically reasonable --- I was not 100\% sure at this point
whether the boundary conditions were correct.

\failure
\begin{itemize}[nosep]
  \item Boundary conditions were wrong --- I initially fixed only the bottom
    face vertices, but this created a cantilever condition rather than a
    foundation condition.
  \item Several structures \textbf{collapsed} in simulation --- the solver
    produced displacements $>1$\,m, which is physically nonsensical for a
    concrete building.
  \item Solver was extremely slow: 45--90 minutes per building at $128^3$
    mesh resolution.
\end{itemize}

\fix
Added fixed supports at foundation nodes (min-$z$ face, i.e.\ the bottom
slab). This is still an idealized BC but prevents the numerical collapse issue. Also added a sanity
check: if $\|u\|_\infty > 1.0$\,m, flag the result as diverged.

\nextstep
Run FEA on all 14,293 wireframes and filter bad results.

% ──────────────────────────────────────────────────────────────
\entry{February 5, 2026}{Mass FEA Processing and Filtering}

\objective
Generate FEA results for all 14,293 wireframes. Identify and remove bad data.

\experiment
Ran \texttt{solve\_asce7\_22\_asd\_sfepy\_ai\_labels.py} on all wireframes.
8 ASCE 7-22 ASD load combinations:
D, D+L, D+0.6W, D+0.7E, 0.6D+0.6W, 0.6D+0.7E, D+0.75L+0.45W,
D+0.75L+0.525E.

Each output: max Von Mises stress, max displacement, global compliance,
principal stresses. Voxelized at $128^3$ with 7 channels (1 occupancy +
6 part one-hot).

\result
I ran all 14,293 simulations in $\sim$10~days on the lab workstation. The total raw data came to $\sim$200\,GB.

\failure
\textbf{3,115 out of 14,293 samples (21.8\%) produced bad FEA results:}
\begin{itemize}[nosep]
  \item 2,041 --- solver diverged (displacement $> 1.0$\,m)
  \item 612 --- near-zero compliance ($< 10^{-6}$\,J), indicating degenerate geometry
  \item 307 --- Von Mises stress $\leq 0$, physically impossible
  \item 155 --- geometry too thin for meshing (walls $< 1$ voxel after discretization)
\end{itemize}

\fix
I built \texttt{filter\_bad\_data.py} to systematically identify and remove all
bad samples, creating \texttt{clean\_manifest.json} with \textbf{11,178 valid
samples}.

I applied a family-aware train/val/test split to prevent data leakage from
similar wireframes: \textbf{8,943 train / 1,121 val / 1,114 test}.

\nbimgcap[0.70\linewidth]{nb_fea_filter.png}{The breakdown of my 14,293
wireframes after FEA filtering. I kept 78.2\% (11,178 samples) as valid
data. The 21.8\% rejection rate was expected given the geometric diversity.}

\obs
Losing 21.8\% of data is significant but expected. Building wireframes have
huge geometric variability --- some are simply not structurally viable.
The important thing is that the remaining 11,178 have valid, physically
meaningful FEA outputs.

\nextstep
Begin training surrogate models.

% ──────────────────────────────────────────────────────────────
\entry{February 8, 2026}{Voxel Resolution Experiment: 32$^3$ vs.\ 64$^3$ vs.\ 128$^3$}

\objective
Determine what voxel resolution I need for the surrogate to learn
structural behavior accurately.

\experiment
I voxelized buildings at $32^3$, $64^3$, and $128^3$. Then I trained a small CNN
(4 conv layers, 2 FC layers) on each resolution for 50 epochs and measured
Spearman $\rho$ for compliance on the validation set.

\result
\begin{center}
\begin{tabular}{lcc}
  \toprule
  \textbf{Resolution} & \textbf{Spearman $\rho$} & \textbf{Train time/epoch} \\
  \midrule
  $32^3$ & 0.61 & 12\,s \\
  $64^3$ & 0.78 & 45\,s \\
  $128^3$ & 0.89 & 210\,s \\
  \bottomrule
\end{tabular}
\end{center}

\nbimgcap[0.70\linewidth]{nb_resolution_ablation.png}{My resolution ablation
experiment. $32^3$ fails to represent thin interior walls; $128^3$ is clearly
necessary for the surrogate to capture structural behaviour.}

\failure
$32^3$ resolution is clearly insufficient --- thin interior walls (1--2 voxels)
disappear entirely, causing the model to misrepresent structural stiffness.
At $32^3$, many buildings are just solid blocks with no internal structure.

\fix
Selected $128^3$ as the operating resolution. The 4.7$\times$ longer training
time is acceptable because inference (which matters for optimization speed)
scales better than training.

\nextstep
Design the full 3D-CNN architecture for $128^3$ inputs.


% ════════════════════════════════════════════════════════════
\section{Phase 3: Surrogate Model Development}
% ════════════════════════════════════════════════════════════

% ──────────────────────────────────────────────────────────────
\entry{February 12, 2026}{CNN Architecture --- v1 (Baseline)}

\objective
Design a 3D convolutional network to predict stress, displacement, and
compliance from a $128^3$ voxel grid.

\experiment
Architecture: 4 strided Conv3D stages ($\text{stride}=2$ each), ReLU
activations, BatchNorm, followed by global average pooling and 3 FC layers
down to 3 scalar outputs. Input: 7 channels (occupancy + part one-hot).
Total: $\sim$5.2M parameters.

\result
I got converged training. After 100~epochs, compliance $\rho = 0.72$. Not bad for a first attempt.

\failure
\begin{itemize}[nosep]
  \item Displacement $\rho = 0.51$ --- very poor. The model cannot capture
    localized displacement concentrations.
  \item Loss landscape was noisy --- training loss oscillated significantly
    between epochs.
  \item ReLU caused many dead neurons in early layers (checked gradient
    statistics --- 15\% of conv1 neurons had zero gradients).
\end{itemize}

\nextstep
Try GELU activation and add residual connections.

% ──────────────────────────────────────────────────────────────
\entry{February 14, 2026}{Architecture v2 --- GELU + SE Blocks}

\objective
Improve model capacity with GELU activation and Squeeze-and-Excitation blocks.

\experiment
Changes from v1:
\begin{itemize}[nosep]
  \item Replaced ReLU with GELU (smooth, no dead neurons)
  \item Added 3 SE-ResBlocks after the conv stages (channel attention,
    reduction ratio 4)
  \item Added stochastic depth (drop path rate 0.1) for regularization
  \item Pre-norm BatchNorm3d ordering
  \item Dual global pooling: avg + max concatenated
  \item Base channels: 64 (up from 32)
  \item Total: $\sim$8.76M parameters
\end{itemize}

\result
Significant improvement:
\begin{center}
\begin{tabular}{lcc}
  \toprule
  \textbf{Target} & \textbf{v1 $\rho$} & \textbf{v2 $\rho$} \\
  \midrule
  Compliance & 0.72 & 0.84 \\
  Displacement & 0.51 & 0.79 \\
  VM Stress & 0.58 & 0.69 \\
  \bottomrule
\end{tabular}
\end{center}

\obs
SE blocks help a lot for displacement (+0.28 $\rho$). GELU eliminated the dead
neuron problem. But stress prediction is still poor --- localized stress
peaks are hard to capture with global pooling.

\nbimgcap[0.72\linewidth]{nb_arch_comparison.png}{My architecture comparison:
v1 (ReLU, no attention) vs.\ v2 (GELU + SE blocks). The SE blocks gave a
huge +0.28 improvement in displacement prediction.}

\nextstep
Address training stability and target normalization.

% ──────────────────────────────────────────────────────────────
\entry{February 16, 2026}{Target Normalization --- The Orders-of-Magnitude Problem}

\objective
Fix the target distribution issue. FEA outputs span enormous ranges.

\experiment
Analyzed target distributions:
\begin{itemize}[nosep]
  \item VM stress: $10^3$ to $10^8$ Pa (5 orders of magnitude)
  \item Displacement: $10^{-6}$ to $10^{-1}$ m (5 orders)
  \item Compliance: $10^{-2}$ to $10^{6}$ J (8 orders!)
\end{itemize}

Tried three normalization approaches:
\begin{enumerate}[nosep]
  \item Raw values (baseline) --- loss dominated by compliance
  \item Standard z-score --- still heavy tails
  \item log1p $\to$ z-score with 2nd/98th percentile winsorization
\end{enumerate}

\result
Option 3 worked dramatically better:
\begin{center}
\begin{tabular}{lccc}
  \toprule
  & \textbf{Raw} & \textbf{Z-score} & \textbf{log1p+winsor} \\
  \midrule
  Compliance $\rho$ & 0.84 & 0.87 & 0.93 \\
  Loss stability & poor & moderate & good \\
  \bottomrule
\end{tabular}
\end{center}

\failure
\sout{Initially tried to predict min\_safety\_factor as a 4th target.}
Analysis showed it was 99.8\% anti-correlated with VM stress --- completely
redundant. Dropped it to reduce confusion for the ensemble.

\fix
Adopted log1p transform $\to$ winsorize at 2nd/98th percentile $\to$ z-score
normalize. This squashes the extreme tails and ensures all three targets have
similar scale. Loss function: Huber (smooth L1) with per-target weights
(VM=1.0, displacement=1.5, compliance=1.0).

\nbimgcap[0.70\linewidth]{nb_normalization.png}{My normalization comparison.
The log1p $+$ winsorize strategy clearly outperforms raw values or plain z-score
by squashing the extreme tails of the FEA output distribution.}

\nextstep
Train the full ensemble model.

% ──────────────────────────────────────────────────────────────
\entry{February 17, 2026}{Ensemble Training --- v2 (5 Members)}

\objective
Train a 5-member deep ensemble for uncertainty-quantified predictions.

\experiment
Five independently-seeded copies of the v2 architecture. Training:
\begin{itemize}[nosep]
  \item AdamW optimizer, $lr = 5 \times 10^{-4}$, weight decay $10^{-4}$
  \item Cosine annealing with warm restarts (period 15 epochs)
  \item EMA (exponential moving average, $\alpha = 0.999$)
  \item Gradient clipping (max norm 1.0)
  \item Data augmentation: random 90$^\circ$ rotations + reflections
  \item Early stopping: patience 30 epochs
  \item Hardware: 4$\times$ NVIDIA GB200 GPUs
\end{itemize}

\result
I got training started and the first ensemble member converged in $\sim$130~epochs.

\failure
\textbf{GB200 Training Crashes (February 18).} Training crashed repeatedly:
\begin{itemize}[nosep]
  \item CUDA out-of-memory when \texttt{num\_workers} exceeded GPU memory
  \item Data loader workers created zombie processes
  \item No error logs --- crashes were silent
\end{itemize}

\fix
\begin{itemize}[nosep]
  \item Capped \texttt{num\_workers} at 4 per GPU
  \item Added data validation (check tensor shapes before forward pass)
  \item Added comprehensive error logging to catch silent crashes
  \item Added gradient checkpointing to reduce memory footprint
\end{itemize}

\nextstep
Complete ensemble training and evaluate.

% ──────────────────────────────────────────────────────────────
\entry{February 19, 2026}{Data Filtering v3 and Retraining}

\objective
Clean the training data more aggressively after examining ensemble disagreement
patterns.

\experiment
Analyzed samples where ensemble members disagreed most. Found that high
disagreement correlated strongly with specific data quality issues. Created
v3 of the training pipeline:
\begin{itemize}[nosep]
  \item Stricter filter: displacement $> 1.0$\,m OR compliance $< 10^{-6}$ OR
    VM $\leq 0$
  \item Removed \texttt{min\_safety\_factor} target (99.8\% anti-correlated
    with VM stress)
  \item Per-target R$^2$ monitoring during training
  \item Weighted loss (displacement weight = 1.5$\times$)
\end{itemize}

\result
After retraining from scratch on clean data (v3):

\nbimgcap[0.80\linewidth]{nb_training_curves.png}{Training loss curves for all
5 ensemble members. All members converge by epoch $\sim$130. The spread between
members becomes my ensemble uncertainty estimate.}

\obs
Compliance $\rho$ improved from 0.87 to 0.93 just by removing bad training data.
I think this is the right takeaway but I am not 100\% sure whether the
improvement is all data quality or partly the other v3 changes (weighted loss,
dropped safety factor target). Hard to fully isolate.

\nextstep
Push training further and validate on held-out test set.

% ──────────────────────────────────────────────────────────────
\entry{February 20, 2026}{Model v14 --- Breakthrough Validation}

\objective
Evaluate the fully-trained v3 ensemble on the 1,114-sample held-out test set.

\experiment
Ran inference on all test samples. Computed ensemble mean $\mu$ and standard
deviation $\sigma$ for each target. Measured Spearman $\rho$, $R^2$, and MAPE.

\result
\textbf{This felt like the breakthrough moment.} (Writing this the same
evening --- want to mark it because I have been doubting the approach for
weeks.)

\begin{center}
\begin{tabular}{lccc}
  \toprule
  \textbf{Target} & \textbf{Spearman $\rho$} & $R^2_{\log}$ & \textbf{MAPE} \\
  \midrule
  VM stress & 0.737 & 0.419 & 37.4\% \\
  Displacement & 0.970 & 0.842 & 10.9\% \\
  Compliance & \textbf{0.948} & 0.814 & 18.5\% \\
  \bottomrule
\end{tabular}
\end{center}

\obs
\begin{itemize}[nosep]
  \item Compliance $\rho = 0.948$ far exceeds my target of $\geq 0.90$.
  \item VM stress MAPE is high (37.4\%) because stress is localized at
    specific points --- hard to predict from global features.
  \item Displacement $\rho = 0.970$ is very strong --- the model has
    learned global stiffness patterns.
  \item The model does not need pointwise accuracy. It needs to correctly
    \emph{rank} designs so sensitivity scores steer optimization correctly.
  \item \textit{Open question at this point: will $\rho = 0.948$ actually be
    good enough for safe optimization, or will errors compound? I do not
    know yet --- that is what the validation phase will answer.}
\end{itemize}

\fix
Adopted model v14 as the production surrogate. Froze architecture and weights.

\nbimgcap[0.80\linewidth]{nb_surrogate_accuracy.png}{Surrogate regression: my
predicted vs.\ actual compliance on 1,114 held-out test designs (log scale).}

\nextstep
Begin building the topology optimization loop.


% ════════════════════════════════════════════════════════════
\section{Phase 4: Topology Optimization Development}
% ════════════════════════════════════════════════════════════

% ──────────────────────────────────────────────────────────────
\entry{February 21, 2026}{CMA-ES Optimizer --- First Attempt (Abandoned)}

\objective
Try using CMA-ES (evolutionary strategy) for topology optimization,
parameterized by wall-segment thicknesses.

\experiment
Parameterized each building as $\sim$29 continuous variables (one thickness
per wall segment). Used CMA-ES to minimize volume subject to structural
constraints (stress, compliance). Each CMA-ES function evaluation queries
the trained ensemble.

\result
I got CMA-ES to converge on simple test cases (3--5 wall segments), removing
$\sim$12\% of material.

\failure
\begin{itemize}[nosep]
  \item With $>$15 wall segments, CMA-ES needed 500+ evaluations to converge.
  \item The wall-segment parameterization was too coarse --- could only
    uniformly thin walls, not create topological changes.
  \item No control over connectivity --- some designs became disconnected.
  \item Maximum material removal was limited to $\sim$15\%.
\end{itemize}

\fix
\sout{Try increasing CMA-ES population size.} Abandoned CMA-ES entirely.
The wall-segment parameterization is fundamentally too restrictive. I need
\textbf{voxel-level control} to achieve meaningful topology optimization.

\nextstep
Design a gradient-based voxel erosion algorithm.

% ──────────────────────────────────────────────────────────────
\entry{February 21, 2026 (evening)}{Sensitivity-Guided Batch Erosion --- v1}

\objective
Replace CMA-ES with a gradient-based approach: compute per-voxel sensitivity
scores and remove the least structurally important voxels.

\experiment
Algorithm:
\begin{enumerate}[nosep]
  \item Compute distance transform from exterior surface
  \item Find candidate voxels deeper than minimum thickness
  \item Backprop through ensemble: $s_i = \frac{\partial}{\partial p_i}
    [f(C) + 0.3 \cdot f(\sigma)]$
  \item Rank by sensitivity (highest = least important)
  \item Remove top-$B$ voxels in a batch
  \item Check constraints via ensemble --- if violated, undo and halve $B$
\end{enumerate}

\result
On my first test (reference house~00472), I removed 28\% of material in just 45~s.
This immediately blew CMA-ES out of the water.

\failure
\textbf{Disconnected structures!} After $\sim$150 iterations, the algorithm
removed load-bearing columns, producing a mesh with 3 disconnected components.
The ``optimized'' design looked like Swiss cheese.

\obs
The sensitivity score correctly identifies structurally unimportant voxels,
but it does not account for \emph{topology} --- removing a low-sensitivity
voxel can disconnect the structure if it was the only bridge between components.

\nextstep
I need a topology preservation constraint.

% ──────────────────────────────────────────────────────────────
\entry{February 22, 2026}{The 26-Connectivity Catastrophe}

\objective
Add connectivity enforcement to prevent disconnected structures during
optimization.

\experiment
First attempt: check that the voxel field remains connected using
\textbf{26-adjacency} (voxels touching at faces, edges, or corners are
considered connected). A voxel is removable only if its deletion does not
increase the number of 26-connected components.

\result
I verified that every optimized voxel field was a single 26-connected component. I thought I had solved the problem.

\failure
\textbf{Except I hadn't.} When I ran Marching Cubes
to convert the voxels to a triangle mesh for STL export, the mesh had
\textbf{hundreds of disconnected floating fragments.}

The issue: 26-connectivity allows voxels to be ``connected'' through diagonal
corners. But Marching Cubes only generates mesh surfaces where voxels share
\emph{faces}. So a voxel touching only at a corner is topologically connected
in voxel space but becomes an isolated island in mesh space.

\nbimgcap[0.75\linewidth]{fig12_stl_comparison.png}{My exported 3D STL:
original structure (top) vs.\ SASTO-PA optimized (bottom). The interior
walls are visibly thinner while the exterior shell is preserved.}

\fix
\textbf{I switched to 6-connectivity for the occupied voxels.} Two
occupied voxels are 6-adjacent iff they share a face. Under this stricter
criterion, a single 6-connected voxel component \emph{guarantees} a single
Marching Cubes mesh component (Kong \& Rosenfeld 1989).

I implemented the \textbf{6-simple-point test}: a voxel is removable only if
deletion does not change the number of 6-connected components. I check this
in $O(1)$ via a 26-neighborhood Euler characteristic lookup table (pre-computed
$2^{26}$ entries, stored as a hash set).

\nextstep
Validate that 6-connectivity actually produces printable meshes.

% ──────────────────────────────────────────────────────────────
\entry{February 22, 2026 (afternoon)}{6-Connectivity Validation}

\objective
Verify that 6-connected topology enforcement produces single-component
watertight meshes after Marching Cubes.

\experiment
I ran the updated SASTO algorithm (with 6-simple-point) on 60 diverse test
designs. For each, I checked:
\begin{enumerate}[nosep]
  \item Number of connected components in the voxel field
  \item Number of connected components in the Marching Cubes mesh
  \item Mesh watertightness (euler number = 2)
\end{enumerate}

\result
\begin{center}
\begin{tabular}{lc}
  \toprule
  \textbf{Metric} & \textbf{Result} \\
  \midrule
  Single-component voxel fields & 60/60 (100\%) \\
  Single-component meshes (direct) & 54/60 (90\%) \\
  Single-component after CC fill pass & 60/60 (100\%) \\
  Zero geometry discards & 60/60 (100\%) \\
  \bottomrule
\end{tabular}
\end{center}

\obs
The 10\% that needed a CC fill pass had very thin ($\leq 1$ voxel) bridges that
Marching Cubes rounded off. A single connected-component fill resolves these
trivially. None of my designs needed to be discarded --- a massive improvement
over the 26-connectivity approach which had catastrophic failures.

\nextstep
Implement part-aware thickness policy.

% ──────────────────────────────────────────────────────────────
\entry{February 23, 2026}{Part-Aware Thickness --- SASTO-PA vs SASTO-U}

\objective
Allow different minimum wall thicknesses for different structural parts.
Interior walls carry negligible vertical load and can be thinned more
aggressively than exterior walls.

\experiment
Two variants:
\begin{itemize}[nosep]
  \item \textbf{SASTO-U} (Uniform): $t_{\min} = 2$ voxels for all parts.
    This is $2 \times 78.125$\,mm $\approx 156$\,mm.
  \item \textbf{SASTO-PA} (Part-Aware): $t_{\min} = 2$ voxels for exterior
    walls, roof, floor; $t_{\min} = 1$ voxel for interior walls ($\approx 78$\,mm).
\end{itemize}

Tested on reference case (Sample 00472):

\result
\begin{center}
\begin{tabular}{lccc}
  \toprule
  \textbf{Variant} & \textbf{Vol.\ Reduction} & \textbf{Max $C/C_0$} \\
  \midrule
  SASTO-U & 34.3\% & 0.98 \\
  SASTO-PA & \textbf{45.0\%} & 1.004 \\
  \bottomrule
\end{tabular}
\end{center}

\obs
Part-aware thickness gives an extra \textbf{10.7 percentage points} of material
removal! The interior walls are thinned from $\sim$200\,mm to $\sim$80\,mm
while the stress increase stays well within limits.

\nbimgcap[0.75\linewidth]{fig13_voxel_parts.png}{Voxel cross-sections at three
heights (z=20, 50, 80) colored by part label. Interior walls (red) are
candidates for aggressive thinning under SASTO-PA.}

\nextstep
Run SASTO-PA on the full test set.

% ──────────────────────────────────────────────────────────────
\entry{February 23, 2026 (evening)}{Trust-Region Adaptive Batch Size}

\objective
Make the batch size adaptation more robust. Currently, when a constraint is
violated the batch is halved, but this is too aggressive early on and too
conservative at the end.

\experiment
I implemented the three-phase strategy:
\begin{itemize}[nosep]
  \item Phase 1 (bulk): Start with $B=200$. On violation, halve $B$. On
    success, keep $B$ if $>$ min. Stop when $B < 10$.
  \item Phase 2 (endgame): Fix $B=5$, stop after 200 evaluations or when
    $B=1$ and no more removable voxels.
  \item Phase 3 (post-process): Pocket fill, spike removal, SDF smoothing.
\end{itemize}

\result
Reference case: $\sim$260 batches, 52,500 voxels removed, 159.5s total.

\nbimgcap[0.80\linewidth]{nb_convergence.png}{Convergence trace during my
SASTO-PA optimization: volume fraction (top) and compliance ratio (bottom)
vs.\ batch number. The vertical dotted line marks the start of Phase~2.}

\obs
The trust-region works well. Early batches remove 200 voxels each (fast bulk
removal). As we approach constraints, batch sizes shrink automatically. Phase 2
endgame squeezes out the last few percent by removing 1--5 voxels at a time.

\nextstep
Run batch optimization on all 1,114 test geometries.


% ════════════════════════════════════════════════════════════
\section{Phase 5: Large-Scale Batch Optimization}
% ════════════════════════════════════════════════════════════

% ──────────────────────────────────────────────────────────────
\entry{February 25, 2026}{Batch Run --- 1,114 Test Geometries}

\objective
Run SASTO-PA optimization on all 1,114 held-out test designs.

\experiment
I used \texttt{run\_batch\_all.py} (892 lines). For each design:
\begin{enumerate}[nosep]
  \item Load voxel grid and part labels
  \item Run SASTO-PA 3-phase optimization
  \item Export optimized voxel field
  \item Log all metrics (volume, stress, compliance, runtime)
\end{enumerate}

I ran this on my RTX A3000 Laptop GPU (the GB200 cluster was for training only).

\result
I processed all 1,114 designs in $\sim$15.5~hours (median 50~s per design) on my laptop GPU.

\failure
Several issues during the batch run:
\begin{itemize}[nosep]
  \item \textbf{Sample 02962}: CUDA out-of-memory. This house was anomalously
    large ($>$160K occupied voxels). Had to add a hard-coded skip with fallback
    to per-model CPU inference.
  \item 353 designs (38.5\%) had $\leq$1\% material removal (``zero-budget''
    designs). These are houses where the surrogate predicted that any removal
    would violate constraints.
  \item 199 designs (17.9\%) had $>$30\% material removal.
\end{itemize}

\nbimgcap[0.80\linewidth]{nb_histogram.png}{Volume reduction I achieved
across 1,114 test geometries. Median reduction: 23.5\%. The wide spread
reflects geometric diversity in the 3DWire dataset.}

\obs
The ``zero-budget'' designs are mostly houses with dense internal load-bearing
walls. The surrogate is likely being overly conservative on these --- possible
calibration issue. But for a safety-first system, false conservatism is
acceptable.

\nextstep
Validate all optimized designs with independent FEA.

% ──────────────────────────────────────────────────────────────
\entry{February 26, 2026}{CUDA OOM and Fallback Strategy}

\objective
Fix the CUDA out-of-memory crash that killed the batch run on Sample 02962.

\experiment
I used \texttt{run\_batch\_all.py} to run the full population optimization.
For each of the 1,121 validation designs I also did a targeted memory
profiling pass to understand the OOM crash better.

\result
I added a two-level fallback:
\begin{enumerate}[nosep]
  \item Primary: run ensemble on GPU with gradient checkpointing
  \item Fallback: if OOM, move each member to CPU, run sequentially
\end{enumerate}

\obs
I found that the full 5-member ensemble forward pass on a
$128^3$ input requires $\sim$4.2\,GB GPU memory, and backprop for sensitivity
doubles that. Large houses with $>$150K occupied voxels exceed the RTX
A3000's 6\,GB limit.

CPU fallback was $\sim$20$\times$ slower (50s $\to$ $\sim$16min per design).
Only 3 out of 1,114 designs needed this. Acceptable.

\nextstep
Continue batch processing with fallback enabled.


% ════════════════════════════════════════════════════════════
\section{Phase 6: Model Calibration \& $k$-Factor Tuning}
% ════════════════════════════════════════════════════════════

% ──────────────────────────────────────────────────────────────
\entry{February 27, 2026}{$k$-Factor Sweep and Pareto Analysis}

\objective
Tune the conservative bound parameter $k$ in
$\hat{y}^+ = \mu + k \cdot \sigma$. Higher $k$ means more safety margin but
less material removal.

\experiment
I swept $k \in \{0.25, 0.50, 0.75, 1.0, 1.5, 2.0, 3.0\}$ on the validation
set (1,121 designs). For each $k$, I recorded the acceptance rate (fraction
of designs passing constraints) and mean volume reduction.

\nbimgcap[0.80\linewidth]{nb_k_factor.png}{The $k$-factor Pareto sweep I
conducted: acceptance rate (teal) and volume reduction (gold) vs.\ $k$.
I chose $k=1.0$ as the operating point.}

\result
\begin{center}
\begin{tabular}{lccc}
  \toprule
  $k$ & Acceptance & Median Reduction & Safety \\
  \midrule
  0.25 & 92\% & 31\% & loose \\
  0.50 & 96\% & 28\% & moderate \\
  \textbf{1.0} & \textbf{>99\%} & \textbf{23\%} & \textbf{good} \\
  2.0 & 100\% & 14\% & strict \\
  3.0 & 100\% & 6\% & very strict \\
  \bottomrule
\end{tabular}
\end{center}

\fix
I selected $k = 1.0$ as my operating point. This gives a $1\sigma$ safety
buffer against surrogate underprediction while preserving meaningful material
savings.

\nextstep
Lock $k=1.0$ and compute conformal coverage certificates.

% ──────────────────────────────────────────────────────────────
\entry{February 27, 2026 (afternoon)}{Conformal Prediction Certificates}

\objective
Compute distribution-free safety guarantees using conformal prediction theory.

\experiment
On the validation set (1,121 designs), I computed non-conformity scores:
$\alpha_i = |y_i - \mu_i| / \sigma_i$ for each target. I sorted them and
computed quantiles to determine what $k$ is needed for various coverage levels.

\result
\begin{itemize}[nosep]
  \item $k$ for 84.1\% compliance coverage: $k = 1.90$ (heavier tails than
    Gaussian --- makes sense for FEA outputs)
  \item $k$ for 84.1\% VM stress coverage: $k = 4.31$ (very heavy tails due
    to localized stress concentrations)
  \item 99\% conformal upper bound on $C_\text{opt}/C_\text{base}$:
    \textbf{0.950} (vs.\ limit 1.15)
  \item Distribution-free bound (Vovk 2005):
    $P(\text{violation}) \leq \frac{1}{n+1} = 0.09\%$
\end{itemize}

\obs
The compliance tails are heavier than Gaussian ($k_\text{84.1\%} = 1.90$ vs.\
$1.0$ for Gaussian). This means the ensemble underestimates uncertainty on about
$16\times$ more designs than a Gaussian would predict. However, at $k = 1.0$,
the overall safety margin is still very large (conformal bound 0.950 vs.\
limit 1.15 = 0.200 margin).

\nextstep
Lock calibration and run independent FEA validation.


% ════════════════════════════════════════════════════════════
\section{Phase 7: Independent FEA Validation}
% ════════════════════════════════════════════════════════════

% ──────────────────────────────────────────────────────────────
\entry{February 28, 2026}{Building Independent Hex8 FEA Validator}

\objective
Build a completely independent FEA solver to verify SASTO optimized designs.
This must be different from the training FEA to avoid circular validation.

\experiment
I wrote \texttt{run\_fea\_validation.py}: uniform hex8 voxel FEA (not
tetrahedral like training data). I used the same boundary conditions and loads.

I chose the AMG-preconditioned conjugate gradient (\texttt{pyamg} library)
as my solver. This is numerically different from the ScipyDirect solver I used
in training data generation.

\result
My validation solver worked correctly. It needed an average of 83~$\pm$~12 CG iterations and 106~$\pm$~30~s per paired solve (original + optimized).

\obs
AMG preconditioning was critical. Without it, plain CG needed 4,000--5,000
iterations (vs.\ 83 with AMG). I spent a lot of time thinking my solver was
diverging --- it was just converging incredibly slowly without preconditioning.

\nextstep
Validate all 1,114 optimized designs.

% ──────────────────────────────────────────────────────────────
\entry{March 1, 2026}{Full Population FEA Validation ($n = 1{,}114$)}

\objective
Run independent hex8 FEA on all 1,114 optimized designs. This is the ``ground
truth'' safety check.

\experiment
For each design, I ran FEA on both the original and the SASTO-PA optimized
geometry, computing the compliance ratio $C_\text{opt}/C_\text{base}$.\\

\result
\textbf{Zero violations across all 1,114 designs I tested.}

\begin{center}
\begin{tabular}{lr}
  \toprule
  \textbf{Metric} & \textbf{Value} \\
  \midrule
  Mean $C_\text{opt}/C_\text{base}$ & $0.631 \pm 0.112$ \\
  Maximum ratio & \textbf{1.004} \\
  Designs $> 1.15$ & \textbf{0/1,114} \\
  99\% conformal bound & 0.950 \\
  Safety margin to limit & 0.146 units \\
  \bottomrule
\end{tabular}
\end{center}

\nbimgcap[0.80\linewidth]{nb_fea_validation.png}{Compliance ratio
$C_\text{opt}/C_\text{base}$ for all 1,114 designs I tested. The dashed red
line is the 1.15 limit. My maximum was 1.004 --- zero violations.}

\obs
The maximum compliance ratio is 1.004 --- the surrogate's conservative bound
$(\mu + k\sigma)$ maintained a massive safety margin. Most optimized designs
are actually \emph{more} compliant than the original (ratio $< 1.0$), likely
because removing sparse interior voxels doesn't significantly affect global
stiffness.

\nextstep
Run SIMP benchmark comparison.

% ──────────────────────────────────────────────────────────────
\entry{March 1, 2026 (afternoon)}{SIMP Benchmark Comparison}

\objective
Compare SASTO runtime against classical SIMP topology optimization at
equivalent resolution.

\experiment
I ran standard SIMP (density-based) optimization on 10 selected test designs at
$64^3$ resolution (not $128^3$ --- SIMP was too slow at that resolution).
I measured per-design runtime and projected $128^3$ times using $O(N^2)$ scaling.

\result
\begin{center}
\begin{tabular}{lcc}
  \toprule
  \textbf{Method} & \textbf{Per Design} & \textbf{vs.\ SASTO} \\
  \midrule
  SASTO (ours) & 50\,s median & baseline \\
  SIMP $64^3$ & 94\,s median & $1.9\times$ slower \\
  SIMP $128^3$ (projected) & 19--77\,min & $23$--$92\times$ slower \\
  \bottomrule
\end{tabular}
\end{center}

\nbimgcap[0.80\linewidth]{nb_speedup.png}{My runtime comparison (log scale).
SASTO achieves 50~s median vs.\ SIMP's 19--77~min projected at $128^3$.}

\obs
SIMP actually achieves slightly higher reduction (31\% vs.\ 28\% median at
$64^3$) but takes 23--92$\times$ longer. For practical design exploration
(trying dozens of variants), SASTO is far more useful.

\nextstep
Generate final figures and validation plots.


% ════════════════════════════════════════════════════════════
\section{Phase 8: Final Validation \& Figure Generation}
% ════════════════════════════════════════════════════════════

% ──────────────────────────────────────────────────────────────
\entry{March 2, 2026}{Ensemble Uncertainty During Optimization}

\objective
Visualize how ensemble uncertainty changes throughout the optimization process
for the reference case.

\experiment
I logged ensemble $\mu$ and $\sigma$ at every iteration during reference case
optimization and plotted $\mu \pm \sigma$ bands.

\nbimgcap[0.80\linewidth]{nb_convergence.png}{Convergence trace for the reference
case. I logged the ensemble $\mu$ and $\sigma$ at each batch; uncertainty
bands narrow as the design approaches a clean, efficient topology.}

\obs
Uncertainty $\sigma$ naturally decreases as optimization converges. This makes
physical sense --- heavily optimized designs with clean geometry are better
represented in the training distribution. The ensemble is effectively
``more confident'' on designs that look like good structures.

This self-calibrating behaviour is a nice emergent property that provides
additional implicit safety precisely when the design is most aggressively
pruned.

% ──────────────────────────────────────────────────────────────
\entry{March 2, 2026 (afternoon)}{Per-Part Volume Retention Analysis}

\objective
Understand which structural components lose the most material during
optimization.

\experiment
I computed per-part volume retention for all 1,114 optimized designs
(exterior walls, interior walls, roof, floor).

\result
\begin{center}
\begin{tabular}{lc}
  \toprule
  \textbf{Part} & \textbf{Mean Retained (\%)} \\
  \midrule
  Exterior walls & 91.6\% \\
  Interior walls & \textbf{45.3\%} \\
  Roof & 96.8\% \\
  Floor & 98.2\% \\
  \bottomrule
\end{tabular}
\end{center}

\nbimgcap[0.80\linewidth]{nb_per_part.png}{Per-part volume retention I measured
across 1,114 test designs. Interior walls are thinned most (45.3\% retained),
while roof and floor are nearly untouched.}

\obs
Interior walls are thinned most aggressively --- makes structural sense since
they carry minimal vertical load. Roof and floor are almost untouched because
they are critical load paths (roof distributes wind/snow, floor provides base
support). The SASTO-PA thickness floor ($t_\text{min} = 1$ voxel for interior)
allows this selective thinning.

% ──────────────────────────────────────────────────────────────
\entry{March 3, 2026}{STL Export and Mesh Quality}

\objective
Export the optimized designs as watertight STL files suitable for 3D printing.

\experiment
Post-processing pipeline:
\begin{enumerate}[nosep]
  \item Pocket fill: flood-fill enclosed voids
  \item SDF smoothing: Gaussian-smoothed distance field, $\sigma = 0.5$ voxel
  \item Marching Cubes at $128^3$
  \item Mesh simplification to $\sim$35K faces
  \item Watertightness check (Euler number = 2, no self-intersections)
\end{enumerate}

\result
All 1,114 designs exported as valid STL files. Mean file size: 1.8\,MB.
90\% were single-component directly; 10\% required one CC fill pass.

\nbimgcap[0.72\linewidth]{nb_connectivity.png}{The connectivity fix I
implemented: 26-adjacency (left) allows corner-touching voxels that break
Marching Cubes meshes, while 6-adjacency (right) guarantees single-component
output.}

% ──────────────────────────────────────────────────────────────
\entry{March 3, 2026 (evening)}{Bland-Altman Analysis}

\objective
Check for systematic bias in surrogate predictions using Bland-Altman plots.

\experiment
For each of the 1,114 test designs, I computed:
\begin{itemize}[nosep]
  \item Mean of surrogate and FEA compliance: $(y_\text{surr} + y_\text{FEA})/2$
  \item Difference: $y_\text{surr} - y_\text{FEA}$
\end{itemize}
I then plotted the difference vs.\ mean with $\pm 1.96$ SD limits.

\nbimgcap[0.80\linewidth]{nb_bland_altman.png}{Bland-Altman plot I generated:
my surrogate vs.\ independent FEA compliance. Most designs fall within
$\pm 1.96$~SD limits and I found no systematic bias.}

\obs
No systematic bias detected. The surrogate slightly overpredicts compliance
for high-compliance designs (conservative direction, which is good for safety).

% ──────────────────────────────────────────────────────────────
\entry{March 5, 2026}{Scaling Law Analysis}

\objective
Estimate how much additional training data would improve surrogate accuracy.

\experiment
I trained the surrogate on random subsets of the training data (25\%, 50\%,
75\%, 100\%) and measured compliance $\rho$ on the validation set. I then fit a
power law: $\rho(N) = a - b \cdot N^{-c}$.

\nbimgcap[0.70\linewidth]{nb_scaling_law.png}{The data scaling law I fitted.
Diminishing returns appear beyond $\sim$15K samples; my current 8,943
samples already place me near the elbow of the curve.}

\obs
Diminishing returns are setting in. Going from 8,943 to 15,000 training samples
would improve $\rho$ from 0.948 to $\sim$0.960. This suggests the current
dataset size is near-optimal for the architecture --- further gains would
require architectural changes (e.g., attention mechanisms, larger model).


% ════════════════════════════════════════════════════════════
\section{Phase 9: Paper, Poster, \& Final Presentation}
% ════════════════════════════════════════════════════════════

% ──────────────────────────────────────────────────────────────
\entry{March 6, 2026}{FEA Stress Visualization (Fig.~16)}

\objective
Create a publication-quality figure showing FEA stress distribution on the
reference case, before and after optimization.

\experiment
I used pyrender to create 3D stress-colored renderings. I applied the jet
colormap mapped to Von Mises stress and added a horizontal colorbar with a
5\,MPa safety threshold marker.

\nbimgcap[0.80\linewidth]{nb_fea_validation.png}{My independent FEA
validation results: compliance ratios across all 1,114 test designs,
showing zero designs exceed the 1.15 safety limit.}

\failure
My first version had the colorbar overlapping the house rendering. I fixed it by
moving the colorbar to horizontal orientation below panel~(b) only. I also had to
change the threshold line from horizontal to vertical to match the new
orientation.

% ──────────────────────────────────────────────────────────────
\entry{March 7, 2026}{Diverse Geometry Gallery}

\objective
Show that SASTO generalizes across very different house archetypes.

\experiment
I selected 4 representative test geometries covering different plan types
(rectangular, L-shaped, multi-gabled, open-plan) and rendered original, optimized
solid, and cutaway views for each.

\result
I saw material reductions ranging from 18\% (geometry 01440, rectangular plan)
to 45\% (geometry 04203, large open-plan interior). All four passed FEA
validation with zero violations.

\begin{center}
\begin{tabular}{lcc}
  \toprule
  \textbf{Design} & \textbf{Plan Type} & \textbf{Reduction} \\
  \midrule
  01440 & Rectangular & 18.6\% \\
  04203 & Open-plan   & 45.0\% \\
  05728 & L-shaped    & 29.8\% \\
  REF   & Multi-gabled & 27.4\% \\
  \bottomrule
\end{tabular}
\end{center}

\obs
The ``compressibility'' of a design is primarily determined by the fraction
of interior partition surface area relative to total structure. Open-plan
houses allow more aggressive thinning.


% ──────────────────────────────────────────────────────────────
\entry{March 7, 2026 (evening)}{SYNOPSYS Poster v5 --- Full Layout}

\objective
Design a 48'' $\times$ 36'' SYNOPSYS poster with professional layout.

\experiment
I built the poster in LaTeX (\texttt{poster\_v5.tex}) with a three-panel layout:
\begin{itemize}[nosep]
  \item Left (12''): Abstract, Introduction, Problem Framing
  \item Center (24''): Methodology, Results, FEA Validation
  \item Right (12''): Statistical Analysis, Conclusions, Future Work
\end{itemize}
Dark navy background, white content cards, gold accent stripe.

\result
The poster compiled successfully. I placed 15 figures and verified all text
was legible at poster distance ($\geq 24$\,pt fonts).

\failure
First poster version (PowerPoint, automated) had distorted images and looked
unprofessional. Second version (PowerPoint v2) had gradient rendering bugs.
Third version (LaTeX) was much cleaner. Required multiple iterations on
figure placement and spacing.

% ──────────────────────────────────────────────────────────────
\entry{March 8, 2026}{Optimization \& Calibration Diagrams}

\objective
Create clear flowchart diagrams for the SASTO optimization loop, model
calibration procedure, and output packaging pipeline.

\experiment
I generated three separate matplotlib diagrams:
\begin{enumerate}[nosep]
  \item \textbf{Optimization Loop}: input $\to$ distance transform $\to$
    sensitivity $\to$ batch select $\to$ ensemble query $\to$ decision
    diamond (commit/undo) with loop-back
  \item \textbf{Model Calibration}: validation set $\to$ non-conformity scores
    $\to$ $k$-factor fitting $\to$ Pareto search $\to$ lock $k=1.0$
  \item \textbf{Output Packaging}: pocket fill $\to$ SDF smoothing $\to$
    marching cubes $\to$ connectivity check $\to$ mesh QC $\to$ STL export
\end{enumerate}

% ──────────────────────────────────────────────────────────────
\entry{March 9, 2026}{Final Results Summary}

\objective
Compile the complete validated results for SYNOPSYS presentation.

\experiment
I compiled all key results across the 1,114 test geometries:

\result
\begin{center}
\begin{tabular}{lr}
  \toprule
  \textbf{Metric} & \textbf{Value} \\
  \midrule
  Dataset (valid) & 11,178 (from 14,293 wireframes) \\
  Train / val / test & 8,943 / 1,121 / 1,114 \\
  Ensemble members & 5 $\times$ 8.76M params \\
  Compliance Spearman $\rho$ & 0.948 \\
  Displacement $\rho$ & 0.970 \\
  Median volume reduction & 23.5\% \\
  Max reduction (ref.\ case) & 45.0\% \\
  FEA violations & \textbf{0/1,114} \\
  Max compliance ratio & 1.004 (limit: 1.15) \\
  Median runtime & 50\,s \\
  Speedup vs.\ SIMP & 23--92\,\texttimes \\
  Single-component meshes & 90\% direct, 100\% after CC fill \\
  \bottomrule
\end{tabular}
\end{center}

\obs
All three primary endpoints are met:
\begin{enumerate}[nosep]
  \item \textbf{Speed}: 50s median (23--92$\times$ vs.\ SIMP). \checkmark
  \item \textbf{Printability}: 100\% single-component after trivial CC fill.
    \checkmark
  \item \textbf{Safety}: 0/1,114 FEA violations; conformal bound 0.950 vs.\
    limit 1.15. \checkmark
\end{enumerate}


% ════════════════════════════════════════════════════════════
\section{Appendix: Software and Parameters}
% ════════════════════════════════════════════════════════════

\subsection{Software Environment}
\begin{tabular}{ll}
  \toprule
  \textbf{Component} & \textbf{Version / Tool} \\
  \midrule
  Language & Python 3.10 \\
  Deep learning & PyTorch 2.1.0 \\
  FEA (training) & SfePy 2024.1 (tetrahedral) \\
  FEA (validation) & Custom hex8 voxel solver + pyamg \\
  Meshing & Gmsh 4.12, trimesh 4.0 \\
  Boolean operations & manifold3d \\
  3D rendering & pyrender, Open3D \\
  Visualization & matplotlib 3.8, PIL \\
  Hardware (training) & 4$\times$ NVIDIA GB200 \\
  Hardware (optimization) & RTX A3000 Laptop GPU \\
  \bottomrule
\end{tabular}

\subsection{Key Hyperparameters}
\begin{tabular}{ll}
  \toprule
  \textbf{Parameter} & \textbf{Value} \\
  \midrule
  Voxel resolution & $128^3$ \\
  Input channels & 7 (1 occupancy + 6 part one-hot) \\
  Base channels & 64 \\
  SE reduction & 4 \\
  Drop path rate & 0.1 \\
  Learning rate & $5 \times 10^{-4}$ \\
  Weight decay & $10^{-4}$ \\
  EMA decay & 0.999 \\
  Loss function & Huber (smooth L1) \\
  Target transform & log1p $\to$ winsorize (2/98\%) $\to$ z-score \\
  Batch size (training) & 16 per GPU \\
  Optimizer & AdamW \\
  Scheduler & CosineAnnealingWarmRestarts (period=15) \\
  Early stopping patience & 30 epochs \\
  \midrule
  $k$-factor (operating) & 1.0 \\
  Max VM stress bound & $5.0 \times 10^6$ Pa \\
  Max compliance ratio & 1.15 \\
  Min thickness (exterior) & 2 voxels ($\approx$ 156 mm) \\
  Min thickness (interior) & 1 voxel ($\approx$ 78 mm) \\
  Initial batch size & 200 voxels \\
  Min batch (Phase 1) & 10 voxels \\
  Endgame batch & 5 $\to$ 1 voxels \\
  \bottomrule
\end{tabular}

\subsection{Dataset Statistics}
\begin{tabular}{ll}
  \toprule
  \textbf{Statistic} & \textbf{Value} \\
  \midrule
  Source dataset & 3DWire (Lin et al.\ 2024) \\
  Total wireframes & 14,293 \\
  Valid after FEA filter & 11,178 (78.2\%) \\
  Rejected (diverged) & 3,115 (21.8\%) \\
  Training set & 8,943 \\
  Validation set & 1,121 \\
  Test set & 1,114 \\
  Split strategy & Family-aware (prevent leakage) \\
  Raw data volume & $\sim$200 GB \\
  Processed data volume & $\sim$59 GB \\
  \bottomrule
\end{tabular}

% ── END ──────────────────────────────────────────────────────
\vfill
\begin{center}
  \color{notegray}\small\itshape
  End of Project Notebook --- 39 entries covering November 2025 to March 2026.\\
  Reconstructed from dated files, commit logs, and contemporaneous notes.\\
  Errors are crossed out rather than deleted.
\end{center}

\end{document}
